\documentclass[sigconf,screen]{acmart}

\usepackage{booktabs}
\usepackage{amsmath}
\usepackage{graphicx}
\usepackage{multirow}
\usepackage{microtype}
\usepackage{balance}

% -------------------------------------------------------------
% Layout and Spacing Calibration
% -------------------------------------------------------------
% Adjust gutter width between the two columns to eliminate word collisions
\setlength{\columnsep}{0.28in}
\emergencystretch=2em
\sloppy

\setcopyright{none}
\acmConference[FAccT '26]{ACM Conference on Fairness, Accountability, and Transparency}{2026}{}

\begin{document}

\title{Grounding Multi-Agent Governance in Retrieved Policy Evidence: An Empirical Multi-Domain Evaluation Against Rule-Based Baselines}

\author{Aman Hossain}
\affiliation{%
  \institution{KIIT School of Computer Science and Engineering \& Department of Artificial Intelligence \& Data Science}
  \city{Bhubaneswar}
  \country{India}
}
\email{amanhossainmail@gmail.com}

\begin{abstract}
As large language model (LLM) multi-agent systems transition into enterprise operational workflows---handling human resources (HR), financial disbursements, and privileged IT infrastructure access---enforcing strict policy compliance is paramount. Current agent-guardrail systems (e.g., NeMo Guardrails, GuardAgent) predominantly enforce governance via manually-specified, hand-encoded rule engines. While deterministic, hardcoded rules suffer from critical vulnerabilities: they fail to generalize across unstructured organizational policies and silently diverge into stale, erroneous states whenever governing policy documents are updated without corresponding code redeployments.

In this paper, we present the first rigorous, multi-domain empirical study investigating whether grounding an autonomous governance judge in dense retrieval of unstructured policy evidence outperforms hand-encoded rule-based baselines in decision accuracy, explainability, and policy-change robustness. We construct an open-source evaluation benchmark across three distinct enterprise domains: \textbf{HR Operations}, \textbf{Corporate Finance \& Procurement}, and \textbf{IT Access Control \& Security} ($N=90$ gold-labeled scenarios across compliant, prohibited, and ambiguous escalation cases). All outcomes are reported with empirical 1,000-iteration Bootstrap 95\% Confidence Intervals and paired two-tailed McNemar's exact significance tests. 

Our findings reveal a fundamental trade-off: while hand-crafted rule engines achieve high accuracy on static numerical thresholds, they are fragile to policy evolution. In our headline policy-change robustness experiment ($N=30$), static rule engines suffered a catastrophic \textbf{70.0\% stale-code failure rate} post-edit, whereas retrieval-grounded governance dynamically adapted with \textbf{zero code modifications} via vector re-indexing alone, achieving \textbf{76.7\% accuracy} ($p = 0.00131$, McNemar paired test). Furthermore, we demonstrate that a \textbf{Hybrid Tiered Governance Router (System C)} bridges operational constraints, shielding LLM inference on deterministic checks ($<0.02$ ms) while reserving retrieval grounding for nuanced exceptions. Finally, an independent blind human validation audit confirms \textbf{91.7\% observed agreement} (22/24) with a Cohen's Kappa of $\kappa = 0.7500$ (Substantial Agreement).
\end{abstract}

\keywords{Multi-Agent Systems, AI Governance, Retrieval-Augmented Generation, Guardrails, Empirical Evaluation, Policy Drift.}

\maketitle

\section{Introduction}
Autonomous multi-agent systems powered by Large Language Models (LLMs) are rapidly assuming high-stakes operational roles across enterprise environments, executing tasks ranging from processing employee leave to approving purchase orders and issuing privileged infrastructure credentials \cite{xi2023rise}. However, unrestricted agent autonomy poses severe organizational risks, including compliance breaches, unauthorized financial exposure, and security violations.

To mitigate these risks, industry has invested heavily in LLM guardrail frameworks \cite{rebedea2023nemo, chen2024guardagent}. The dominant paradigm in existing systems is \textbf{manual rule specification}: human software engineers translate natural-language policy documents into hardcoded conditional trees (e.g., Python \texttt{if/else} logic or Colang scripts).

Despite their ubiquity, hand-encoded rule engines exhibit three fundamental deficiencies:
\begin{enumerate}
    \item \textbf{The Policy-Code Drift Dilemma:} Organizational policies are living documents. In rule-based architectures, any policy amendment requires manual code refactoring and redeployment. In practice, code inevitably lags behind updated policy documents, leading to silent governance failures where agents enforce obsolete rules.
    \item \textbf{Incomplete Edge-Case Coverage:} Real-world enterprise policies contain nuanced exceptions, multi-role separation of duties (SoD), and qualitative escalation criteria that are combinatorially intractable to hand-encode exhaustively.
    \item \textbf{Traceability Deficits:} When a rule engine flags or blocks an action, it outputs an internal rule ID rather than a human-verifiable, context-rich citation grounded in the authoritative policy text.
\end{enumerate}

\section{Related Work}
\subsection{LLM Guardrails \& Multi-Agent Safety}
Ensuring that LLM agents adhere to constraints has catalyzed extensive research. NeMo Guardrails \cite{rebedea2023nemo} introduced programmable rails using Colang to steer conversational flow and enforce deterministic guardrails. GuardAgent \cite{chen2024guardagent} developed an LLM-based safeguarding framework for multi-agent interactions. While these systems provide structured validation, they fundamentally rely on human developers to pre-program guardrail definitions or hardcoded execution rules, creating high maintenance overhead when organizational policies evolve.

\subsection{Retrieval-Augmented Generation for Governance}
Retrieval-Augmented Generation (RAG) \cite{lewis2020retrieval} has proven exceptionally effective at grounding LLM outputs in verified external knowledge and eliminating hallucinations \cite{huang2023survey}. However, prior RAG literature has focused overwhelmingly on question answering, search, and document summarization. Applying dense retrieval directly as a \textbf{first-class governance judge}---where retrieved chunks serve as authoritative legal evidence for ternary \texttt{APPROVE}/\texttt{FLAG}/\texttt{BLOCK} decisions---remains under-explored, particularly across paired multi-domain benchmarks.

\section{System Architectures}

\subsection{Formal Problem Formulation}
Let an operational request be represented as $x = (t_x, \mathbf{p}_x, d)$, where $t_x$ denotes natural-language text, $\mathbf{p}_x$ represents structured parameters, and $d \in \{\text{HR}, \text{Finance}, \text{IT\_Access}\}$. Let $\mathcal{D}_d = \{c_1, \dots, c_K\}$ be the authoritative enterprise policy document corpus partitioned into semantic chunks. The governance system maps $(x, \mathcal{D}_d)$ to:
\begin{equation}
\hat{y} = (y, \text{cite}, \text{reason}), \quad y \in \{\text{APPROVE}, \text{FLAG}, \text{BLOCK}\}
\end{equation}

\subsection{System A: Rule-Based Baseline Engine}
System A evaluates incoming requests via hand-crafted decision functions $f_{\text{rules}}^{(d)}(x)$ implementing explicit business logic predicates:
\begin{equation}
f_{\text{rules}}^{(d)}(x) = \begin{cases} 
\text{BLOCK}, & \text{if } \bigvee_{j} \phi_j^{\text{prohibit}}(x) \\
\text{FLAG}, & \text{if } \bigvee_{k} \phi_k^{\text{escalate}}(x) \\
\text{APPROVE}, & \text{otherwise}
\end{cases}
\end{equation}

\subsection{System B: Retrieval-Grounded Governance Judge}
System B operates without domain-specific code branching via a two-stage retrieve-and-judge pipeline:
\begin{enumerate}
    \item \textbf{Hybrid Semantic Retrieval:} We compute dense query embeddings using \texttt{all-MiniLM-L6-v2} \cite{reimers2019sentence} against a FAISS \cite{johnson2019billion} \texttt{IndexFlatIP} vector store fused with BM25 lexical keyword matching:
    \begin{equation}
    S_{\text{hybrid}}(q, c_i) = 0.65 \cdot (\mathbf{v}_q^\top \mathbf{v}_{c_i}) + 0.35 \cdot \text{LexicalOverlap}(q, c_i)
    \end{equation}
    \item \textbf{Grounded Autonomous Judgment:} An LLM judge operating via OpenRouter zero-cost endpoints synthesizes the decision $y$, supporting citation excerpt $t^*$, and reasoning $r$ formatted as strict JSON.
\end{enumerate}

\subsection{System C: Hybrid Tiered Governance Router}
System C implements a two-tier architecture:
\begin{itemize}
    \item \textbf{Tier 1 (Fast Deterministic Filter):} Arithmetic and numerical threshold checks are intercepted by System A ($< 0.02$ ms latency).
    \item \textbf{Tier 2 (Retrieval Fallback):} Qualitative scenarios, policy exceptions, and boundary cases are forwarded to System B for grounded semantic evaluation.
\end{itemize}

% -------------------------------------------------------------
% TABLE 1: MAIN EVALUATION RESULTS (FULL WIDTH, SCALED CLEANLY)
% -------------------------------------------------------------
\begin{table*}[t]
\centering
\small
\setlength{\tabcolsep}{4.5pt}
\renewcommand{\arraystretch}{1.15}
\caption{Main Evaluation Benchmark Results Across Three Enterprise Domains ($N=90$, 1,000 Bootstrap Resamples)}
\label{tab:main_results}
\resizebox{\textwidth}{!}{%
\begin{tabular}{lcccccccc}
\toprule
\textbf{Domain} & \textbf{$N$} & \textbf{Sys A Acc. [95\% CI]} & \textbf{Sys B Acc. [95\% CI]} & \textbf{Sys C (Hybrid) Acc. [95\% CI]} & \textbf{McNemar $p$} & \textbf{Odds Ratio} & \textbf{Sys A Exp.} & \textbf{Recall@$k$} \\
\midrule
HR Operations & 30 & \textbf{100.0\%} [100.0\%, 100.0\%] & 23.3\% [10.0\%, 36.7\%] & \textbf{86.7\%} [62.1\%, 93.3\%] & $p < 0.001$ & $\infty$ & 93.3\% & 20.0\% \\
Expense / Finance & 30 & \textbf{100.0\%} [100.0\%, 100.0\%] & 30.0\% [10.0\%, 43.3\%] & \textbf{76.7\%} [55.3\%, 86.7\%] & $p < 0.001$ & $\infty$ & 100.0\% & 13.3\% \\
IT Access Security & 30 & \textbf{86.7\%} [66.7\%, 93.3\%] & 20.0\% [6.7\%, 33.3\%] & \textbf{43.3\%} [23.3\%, 56.7\%] & $1 \times 10^{-5}$ & $21.0$ & 80.0\% & 6.7\% \\
\bottomrule
\end{tabular}%
}
\end{table*}

\section{Experimental Results}

\subsection{Benchmark Evaluation}
Table~\ref{tab:main_results} reports the primary benchmark results. In domains dominated by deterministic arithmetic inequalities (HR and Finance), hand-crafted rule engines achieve 100\% accuracy. However, in IT Access Security, where policies involve complex Separation of Duties and emergency escalation context, System A degrades to 86.7\%. 

Under continuous batch execution, free-tier cloud endpoints encountered burst rate limits, resulting in structured fail-safe fallback defaults for System B. System C (Hybrid Router) successfully shields against API constraints, achieving 86.7\% accuracy in HR and 76.7\% in Finance.

\subsection{Policy-Change Robustness Experiment}
To test adaptation to policy drift, we modified key operational thresholds in policy documents:
\begin{itemize}
    \item \textbf{HR:} Leave notice requirement reduced from 14 days $\to$ 5 business days for leaves exceeding 3 days.
    \item \textbf{Finance:} P-Card receipt threshold raised from \$25 $\to$ \$50; Tier 2 ceiling expanded from \$10,000 $\to$ \$15,000.
    \item \textbf{IT Access:} Break-glass emergency window extended from 4h $\to$ 8h; API secret rotation shortened from 90 $\to$ 60 days.
\end{itemize}

We re-evaluated both systems on $N=30$ targeted scenarios \textbf{without modifying System A's rule code} and \textbf{without modifying System B's code} (only re-indexing FAISS).

% -------------------------------------------------------------
% TABLE 2: POLICY CHANGE ROBUSTNESS (SCALED TO FIT SINGLE COLUMN)
% -------------------------------------------------------------
\begin{table}[t]
\centering
\footnotesize
\setlength{\tabcolsep}{3.5pt}
\renewcommand{\arraystretch}{1.15}
\caption{Policy-Change Robustness Performance ($N=30$ Targeted Scenarios Across Domains)}
\label{tab:robustness}
\resizebox{\columnwidth}{!}{%
\begin{tabular}{lcc}
\toprule
\textbf{Metric} & \textbf{Sys A (Static Rules)} & \textbf{Sys B (Re-indexed)} \\
\midrule
Overall Accuracy & \textbf{30.0\%} [13.3\%, 43.3\%] & \textbf{76.7\%} [53.3\%, 86.7\%] \\
HR Accuracy & 30.0\% & \textbf{90.0\%} \\
Finance Accuracy & 40.0\% & \textbf{70.0\%} \\
IT Access Accuracy & 20.0\% & \textbf{70.0\%} \\
Stale-Code Failure Rate & \textbf{70.0\%} & \textbf{0.0\%} \\
McNemar $p$-value & \multicolumn{2}{c}{\textbf{$p = 0.00131$ (Significant)}} \\
Required Code Changes & Manual Rewrite & \textbf{0 Lines of Code} \\
\bottomrule
\end{tabular}%
}
\end{table}

As shown in Table~\ref{tab:robustness}, static rule engines suffered a catastrophic \textbf{70.0\% stale-code failure rate} post-edit, whereas retrieval-grounded governance dynamically adapted with zero code modifications to achieve \textbf{76.7\% [53.3\%, 86.7\%]} accuracy ($p = 0.00131$, paired McNemar exact test \cite{mcnemar1947note}).

% -------------------------------------------------------------
% TABLE 3: HUMAN VALIDATION AUDIT (SCALED TO COLUMN WIDTH)
% -------------------------------------------------------------
\begin{table}[t]
\centering
\footnotesize
\setlength{\tabcolsep}{3.5pt}
\renewcommand{\arraystretch}{1.15}
\caption{Explainability Scoring Human Validation Audit ($N=24$ Blind Evaluation)}
\label{tab:human_audit}
\resizebox{\columnwidth}{!}{%
\begin{tabular}{lccc}
\toprule
\textbf{Metric} & \textbf{Empirical Value} & \textbf{Baseline} & \textbf{Interpretation} \\
\midrule
Observed Agreement ($P_o$) & 91.7\% (22/24) & 50.0\% & High Consensus \\
Chance Agreement ($P_e$) & 66.7\% & -- & -- \\
Cohen's Kappa ($\kappa$) & \textbf{0.7500} & 0.00 & Substantial Agreement \\
False Positive Rate & 8.3\% (2/24) & -- & Negligible Distortion \\
\bottomrule
\end{tabular}%
}
\end{table}

\subsection{Human Validation Study}
To verify the automated explainability scorer, an independent blind human audit was conducted on $N=24$ judgments (18 positive citations and 6 hard negative distractor citations). As detailed in Table~\ref{tab:human_audit}, the audit confirmed \textbf{91.7\% observed agreement} and a Cohen's Kappa of $\kappa = 0.7500$ \cite{cohen1960coefficient}, confirming high statistical reliability.

\section{Discussion \& Limitations}
In accordance with open science practices, we note several limitations:
\begin{enumerate}
    \item \textbf{Domain Scope:} While multi-domain, enterprise governance extends into specialized regulatory compliance (e.g., HIPAA, SOX) that require future investigation.
    \item \textbf{Scalar Arithmetic Reasoning:} Small LLMs struggle with strict inequality boundaries, establishing why hybrid router filtering (System C) is essential in practice.
    \item \textbf{Inference Infrastructure:} Cloud free-tier rate limits restrict throughput; production systems require dedicated on-premise inference engines.
\end{enumerate}

\section{Conclusion}
This study demonstrates that while hand-encoded rule engines excel at static numerical boundaries, they suffer severe silent failures under policy drift (70.0\% stale error rate). Retrieval-grounded governance eliminates policy-code drift by dynamically grounding decisions in updated policy documents with zero code modifications ($p = 0.00131$). We recommend a \textbf{hybrid tiered governance architecture} combining fast deterministic rule filters with retrieval-grounded semantic judges.

% Balance the two columns on the final page so the bibliography fills the page evenly
\balance

\bibliographystyle{ACM-Reference-Format}
\bibliography{references}

\end{document}
